\section*{ЗАКЛЮЧЕНИЕ}
\addcontentsline{toc}{section}{ЗАКЛЮЧЕНИЕ}

В ходе выполнения данной работы был создан программный комплекс с визуальной средой проектирования для создания компьютерных игр. Комплекс включает в себя игровой движок на языке C\# с использованием графического API OpenGL и редактор сцен на платформе Windows Forms.

Разработанный программный комплекс позволяет загружать трёхмерные модели в форматах STL и GLB, управлять их трансформацией (позиция, поворот, масштаб), настраивать параметры камеры и источников света (точечные, направленные, прожектор), а также сохранять и загружать сцены в JSON-формат. Реализованная графическая подсистема поддерживает два типа проекций (перспективную и ортографическую), две модели освещения (Гуро и Блинна-Фонга), а также динамическую компиляцию шейдерных программ на основе характеристик материала. Редактор сцен обеспечивает визуальное управление игровыми объектами, их тегирование и настройку видимости в игровом режиме. Демонстрационное приложение, созданное на разработанном движке, подтверждает его работоспособность и возможность создания простых 3D игр.

Основные результаты работы:

\begin{enumerate}
	\item Проведён анализ предметной области разработки игровых движков и существующих решений (Unity, Unreal Engine, Godot). Выявлены особенности традиционного подхода к разработке игр и обоснована необходимость создания собственного программного комплекса с визуальной средой проектирования.
	
	\item Разработана концептуальная модель игрового движка, построены диаграммы компонентов, последовательности и классов. Определены функциональные и нефункциональные требования к системе.
	
	\item Осуществлено проектирование архитектуры программного комплекса (игровой движок на C\# с графическим API OpenGL и редактор сцен на Windows Forms). Разработаны структура компонентов, пользовательский интерфейс редактора и программный интерфейс движка.
	
	\item Реализована и протестирована программная система. Проведено тестирование в соответствии со сценариями использования (импорт моделей, настройка сцены, рендеринг, сохранение и загрузка), подтвердившее работоспособность всех функций.
\end{enumerate}

Все требования, объявленные в техническом задании, были полностью реализованы, все задачи, поставленные в начале разработки проекта, решены. Разработанный программный комплекс готов к использованию в учебных и демонстрационных целях, а также может служить основой для дальнейшего развития и создания более сложных игровых проектов.  
